% 文献综述与理论假说
\section{文献综述与理论假说}

\subsection{技术变革与就业：理论脉络}

技术变革对就业的影响是劳动经济学的经典议题。Acemoglu 和 Restrepo（2018）构建的任务框架将自动化建模为资本在生产任务中替代劳动的过程，同时引入新任务作为劳动需求的补偿力量。该框架的核心洞见在于：当自动化速度超过新任务创造速度时，劳动在国民收入中的份额下降；反之亦然。Acemoglu 和 Restrepo（2019）进一步区分了替代效应和复职效应，认为教育和培训政策可以增强复职效应，缓解自动化对劳动者的冲击。Autor（2015）从历史经验出发，论证技术变革一贯改变工作内容而非消灭工作总量，列举了 ATM 机与银行柜员、电子表格与会计职业等经典反例，这一乐观视角与"技术性失业"的焦虑形成有益对话。

在实证层面，Acemoglu 和 Restrepo（2020）利用美国劳动力市场数据，发现每千名工人增加一台工业机器人将使本地就业率下降约 0.2 个百分点、工资下降 0.4\%。Acemoglu 和 Restrepo（2022）将任务框架扩展到工资不平等分析，发现自动化对低技能工人的替代效应更为显著，加剧了工资极化。Goos 和 Manning（2007）提出的就业极化假说——高技能和低技能岗位增加、中间技能岗位收缩——已被多国数据验证。Frey 和 Osborne（2017）从职业特征角度估算了美国 702 种职业被计算机化的概率，发现约 47\% 的岗位面临高风险。Goldsmith-Pinkham 等（2020）对任务框架下的实证方法进行了系统综述。

在中国语境下，Zhang 等（2022）基于省级面板数据发现数字经济导致了中国就业极化，高技能和低技能岗位增加，中间技能岗位减少。刘军和陈维涛（2021）使用中国家庭追踪调查（CFPS）数据发现数字经济显著改善了就业质量，且高学历和东部地区群体受益更多。

\subsection{数字经济测度方法}

数字经济的测度是实证研究的起点，目前主要有两种进路。第一种是综合指数法：赵涛等（2020）使用熵值法从互联网普及率、从业人员、移动电话普及率和数字普惠金融指数四个维度构建地级市数字经济综合指数，该方法在中国经济学界被广泛引用。Tian 等（2023）基于省级面板数据，从互联网、电信、软件业和数字金融四个维度合成综合指数。He 等（2024）专门讨论了数字经济指标体系的构建方法和省际差异的收敛性。第二种是单一代理指标法：北京大学数字金融研究中心编制的数字普惠金融指数（郭峰等，2020）覆盖 2011--2021 年省、市、县三级，是数字经济实证研究中使用最广泛的代理变量。本文同时使用两种进路，以综合指数为主要解释变量、数字金融指数为稳健性替换变量，直接比较两种测度方式对实证结论的影响。

\subsection{中国数字经济就业效应的实证证据}

近年来，数字经济对中国就业影响的实证研究快速增长，可按研究主题分为三类。

第一类关注就业规模。王芳（2023）基于 2012--2020 年 30 个省份面板数据，采用固定效应和中介效应模型，发现数字经济显著促进就业规模扩张，经济增长和产业结构优化是主要的中介路径。王伟和李华（2025）使用 2011--2022 年省级面板，进一步验证了数字经济对总就业和城镇就业的双重促进效应。

第二类关注就业结构。赵向豪等（2024）使用 2011--2021 年省级面板，验证了数字经济对就业结构的优化效应——数字经济发展推动劳动力从第一、二产业向第三产业转移。杜传忠和疏爽（2024）从技能偏向型技术进步视角，发现数字化赋能降低了低技能劳动力工资、提高了高技能劳动力工资，间接推动就业技能结构优化。Chen 等（2024）使用中国企业微观数据，发现人工智能应用显著改变了企业内部技能需求结构。

第三类关注就业质量和非线性效应。张颖和辛爱洪（2024）基于 2013--2020 年省级面板，采用面板门槛模型发现数字经济对高质量就业的促进作用呈边际效应递增特征，存在人力资本门槛——这是方法上最接近本文的现有文献。李明和张强（2023）验证了产业结构升级在数字经济与就业质量之间的中介和门槛作用。

\subsection{研究空白与本文定位}

综合以上三支文献，本文在以下方面定位自己的边际贡献。第一，数据窗口：多数已有研究截至 2020 年，本文覆盖 2011--2023 年，将 COVID-19 后数字经济加速期纳入分析。第二，非线性视角：已有门槛研究（张颖和辛爱洪，2024）以人力资本为门槛变量，本文以城镇化为门槛变量，揭示数字经济就业效应的城镇化条件。第三，双指标比较：已有研究通常只使用一种数字经济测度，本文直接比较综合指数与数字金融指数的就业效应差异，为"数字化的哪个维度更重要"提供经验证据。

\subsection{理论假说}

基于上述理论逻辑和文献证据，本文提出以下三个递进的研究假说：

\textbf{假说 1（H1）：}数字经济发展推动省际就业结构从第一、二产业向第三产业转移。数字经济综合指数与第三产业就业占比正相关。其机制在于：数字技术降低了服务业的交易成本和生产门槛，催生了新的服务业岗位，同时传统制造业的数字化改造减少了低端制造岗位需求，净效应为劳动力向第三产业流动。

\textbf{假说 2（H2）：}数字经济对就业结构的促进效应存在城镇化门槛。高城镇化地区拥有更完善的数字基础设施、更高的人力资本密度和更大的市场规模，数字经济的就业结构转换效应在这些地区更强。城镇化水平存在一个临界值，超过该临界值后，数字经济对三产就业的拉动效应显著增强。

\textbf{假说 3（H3）：}数字经济的就业结构效应存在区域异质性。根据技术扩散理论，当一项技术达到成熟阶段后，其边际效应在技术基础相对薄弱但正在追赶的地区最大。在中国语境下，东部地区的数字经济已相对成熟，边际效应递减；中部地区正处于数字化的加速追赶期，可能表现出最大的边际效应；西部地区数字基础设施尚不充分，效应可能尚未显现。
